\hnheader[Wie gut ist ein Regressionsmodell wirklich? Fehlermaße und Varianzerklärung.][MSE $\leftrightarrow$ Gauß-Rauschen, MAE $\leftrightarrow$ Laplace-Rauschen]{Regressions-}{metriken}{MAE, MSE, RMSE, $R^2$ und Adjusted $R^2$}
\hnsrc{hand-notes/ML Notes.pdf (Regression Metrics); Rechenbeispiel dort korrigiert ($R^2=0.874$, nicht $0.884$: $\mathrm{SST}=12.30$); Bezug zu MLE aus dem Kursstoff}

\begin{hngrid}
%
\begin{hnp}[raster multicolumn=2,acc=hnorange]{1}{Notation}
$y_i$ wahr, $\hat y_i$ vorhergesagt, $\bar y$ Mittel der wahren Werte, $n$ Beispiele, $p$ Prädiktoren, Residuum $e_i=y_i-\hat y_i$.\\
Ziel: kleine Fehler \emph{und} Generalisierung (auf dem Testset messen!).
\end{hnp}
%
\begin{hnp}[raster multicolumn=2,acc=hnpurple]{2}{Fehlerbasierte Metriken}
\begin{align*}
\mathrm{MAE}&=\tfrac1n\textstyle\sum|y_i-\hat y_i|\\
\mathrm{MSE}&=\tfrac1n\textstyle\sum(y_i-\hat y_i)^2\\
\mathrm{RMSE}&=\sqrt{\mathrm{MSE}}
\end{align*}
Kleiner ist besser. RMSE hat die Einheit von $y$; MSE/RMSE bestrafen große Fehler stärker.
\end{hnp}
%
\begin{hnp}[raster multicolumn=2,acc=hnteal]{3}{Skizze: gut vs.\ schlecht}
\centering
\begin{tikzpicture}[scale=.8]
  \begin{scope}
    \draw[hnnavy] (0,0) rectangle (1.9,1.9); \draw[gray,dashed] (0,0)--(1.9,1.9);
    \foreach \x/\y in {.3/.35,.6/.55,.9/.95,1.2/1.15,1.5/1.6,1.7/1.75}{\fill[hngreen] (\x,\y) circle(.05);}
    \node[font=\tiny] at (.95,-.25) {gut: nahe an $y=\hat y$};
  \end{scope}
  \begin{scope}[shift={(2.6,0)}]
    \draw[hnnavy] (0,0) rectangle (1.9,1.9); \draw[gray,dashed] (0,0)--(1.9,1.9);
    \foreach \x/\y in {.3/1.4,.6/.3,.9/1.7,1.2/.5,1.5/1.2,1.7/.2}{\fill[hnred] (\x,\y) circle(.05);}
    \node[font=\tiny] at (.95,-.25) {schlecht: weit weg};
  \end{scope}
\end{tikzpicture}
\end{hnp}
%
\begin{hnp}[raster multicolumn=3,acc=hnnavy]{4}{$R^2$ und Adjusted $R^2$}
\[ R^2=1-\frac{\mathrm{SSE}}{\mathrm{SST}}=1-\frac{\sum(y_i-\hat y_i)^2}{\sum(y_i-\bar y)^2} \]
$R^2=1$: perfekt; $R^2=0$: nicht besser als der Mittelwert; $R^2<0$: schlechter als der Mittelwert.
\[ R^2_{\mathrm{adj}}=1-(1-R^2)\frac{n-1}{n-p-1} \]
Adjusted bestraft zusätzliche Prädiktoren $\Rightarrow$ fairer Vergleich bei verschiedener Featureanzahl.
\end{hnp}
%
\begin{hnp}[raster multicolumn=3,acc=hngreen]{5}{Welche Metrik wann?}
\renewcommand{\arraystretch}{1.25}
\begin{tabular}{@{}p{1.9cm}p{6.4cm}@{}}
\hline
\textbf{Metrik}&\textbf{Einsatz}\\\hline
MAE&robust gegenüber Ausreißern, leicht interpretierbar\\
MSE / RMSE&große Fehler sind besonders schlimm\\
$R^2$&Anteil erklärter Varianz, skalenfrei\\
Adj.\ $R^2$&Modelle mit verschiedener Featurezahl vergleichen\\\hline
\end{tabular}
\end{hnp}
%
\begin{hnex}[raster multicolumn=3]{Beispiel: 5 Hauspreise (in 100\,k\$)}
$y=(5.0,6.0,7.5,8.0,9.5)$, $\hat y=(4.6,5.5,8.0,7.2,10.0)$, Fehler $(0.4,0.5,-0.5,0.8,-0.5)$.\\
\hnsol\ MAE $=\frac{0.4+0.5+0.5+0.8+0.5}{5}=\ora{0.54}$; MSE $=\frac{1.55}{5}=0.31$; RMSE $=0.557$.\\
$\bar y=7.2$, SST $=4.84+1.44+0.09+0.64+5.29=12.30$, SSE $=1.55$:\\
$R^2=1-\frac{1.55}{12.30}=\ora{0.874}$ \hlm{(maschinell nachgerechnet)}.
\end{hnex}
%
\begin{hnp}[raster multicolumn=3,acc=hnrose]{6}{Warum genau diese Fehler? (MLE)}
Gauß-Rauschen $\Rightarrow$ MLE $=$ \emph{MSE} minimieren (Least Squares). Laplace-Rauschen $\Rightarrow$ \emph{MAE} (robust). \textbf{Huber-Verlust}: quadratisch für kleine, linear für große Fehler -- Kompromiss. Metrik und Trainingsverlust sollten zum Anwendungsziel passen (\ora{Kapitel Fehleranalyse}: Zielfunktion vs.\ Metrik).
\end{hnp}
%
\begin{hnp}[raster multicolumn=2,acc=hnred]{7}{Fallstricke}
\begin{hncon}
\item nur eine Metrik betrachten
\item $R^2$ auf Trainingsdaten (steigt immer)
\item verschiedene Skalen vergleichen (MSE)
\end{hncon}
\end{hnp}
%
\begin{hnp}[raster multicolumn=2,acc=hngreen]{8}{Weitere Metriken (Web)}
\begin{itemize}
\item MAPE (prozentual, Problem bei $y{\approx}0$)
\item Median AE (sehr robust)
\item Quantil-/Pinball-Verlust
\end{itemize}
\end{hnp}
%
\begin{hnp}[raster multicolumn=2,acc=hnorange]{9}{Key Takeaways}
\begin{hnchk}
\item MAE, MSE, RMSE, $R^2$
\item $R^2=1-\mathrm{SSE}/\mathrm{SST}$
\item Metrik zum Ziel passend wählen
\end{hnchk}
\end{hnp}
%
\begin{hnp}[raster multicolumn=6,acc=hnpurple,colback=hnlilac!30]{?}{Selbsttest: kannst du das erklären?}
\hnquiz{Was bedeutet $R^2=0$?}{Das Modell ist nicht besser als die Vorhersage des Mittelwerts $\bar y$.}{Wann MAE statt MSE?}{Bei Ausreißern; MSE quadriert große Fehler.}{Warum Adjusted $R^2$?}{$R^2$ steigt mit jedem Feature; Adjusted bestraft überflüssige Prädiktoren.}
\end{hnp}
%
\begin{hnbanner}[raster multicolumn=6]
„Keine einzelne Zahl erzählt die ganze Geschichte -- lies mehrere Metriken zusammen.“
\end{hnbanner}
\end{hngrid}
